\documentclass{article}
\usepackage{amsmath,amssymb}
\usepackage{xurl}
\usepackage[hidelinks]{hyperref}
% Shared commands for notes in docs/paper-gaps/.

\DeclareUnicodeCharacter{2080}{\ifmmode{}_{0}\else\textsubscript{0}\fi}
\DeclareUnicodeCharacter{2081}{\ifmmode{}_{1}\else\textsubscript{1}\fi}
\DeclareUnicodeCharacter{2082}{\ifmmode{}_{2}\else\textsubscript{2}\fi}
\DeclareUnicodeCharacter{2099}{\ifmmode{}_{n}\else\textsubscript{n}\fi}

\newcommand{\Lean}{\textsc{Lean}}
\ExplSyntaxOn
\NewDocumentCommand{\leanid}{m}{
  \ifmmode
    \textsf{\detokenize{#1}}
  \else
    \group_begin:
    \urlstyle{sf}
    \tl_set:Nn \l_tmpa_tl {#1}
    \tl_replace_all:Nnn \l_tmpa_tl {\_} {_}
    \exp_args:NV \path \l_tmpa_tl
    \group_end:
  \fi
}
\ExplSyntaxOff
\providecommand{\tr}{\operatorname{tr}}
\newcommand{\tnleanIssueBase}{https://github.com/LionSR/TNLean/issues/}
\newcommand{\tnleanPrBase}{https://github.com/LionSR/TNLean/pull/}
\newcommand{\tnleanDocBase}{https://sirui-lu.com/TNLean/docs/}
\newcommand{\ghissue}[1]{\href{\tnleanIssueBase#1}{GitHub issue \##1}}
\newcommand{\ghpr}[1]{\href{\tnleanPrBase#1}{PR \##1}}
\newcommand{\doclink}[1]{\href{\tnleanDocBase#1.html}{\path{#1}}}

% Dirac notation and the identity operator, as in blueprint/src/macros/common.tex.
% \providecommand keeps note-local definitions authoritative.
\usepackage{dsfont}
\providecommand{\ket}[1]{|#1\rangle}
\providecommand{\bra}[1]{\langle #1|}
\providecommand{\braket}[2]{\langle #1 | #2 \rangle}
\providecommand{\ketbra}[2]{|#1\rangle\!\langle #2|}
\providecommand{\Id}{\mathds{1}}
\providecommand{\one}{\mathds{1}}
\providecommand{\C}{\mathbb{C}}
\providecommand{\MN}[1]{M_{#1}(\C)}
\providecommand{\MD}{M_D(\C)}

% Verdict marker: \gapnote{<kind>}{<status>}. Kind is the policy
% classification (clarification, local-correction, scope-restriction,
% unfaithful, false-source, open-gap); status is open, wip (elimination
% underway), resolved, or historical. Machine-read by the site index;
% prints a verdict line.
\newcommand{\gapnote}[2]{\par\noindent\textbf{Verdict:} #1 (#2).\par}

\title{Finite-Group Commutants: The Unital Qualification}
\author{}
\date{2026-10-02}
\begin{document}
\maketitle
\gapnote{local-correction}{resolved}

\section{The source and its necessary qualification}
Schuch, Cirac, and P\'erez-Garc\'ia state that every complex matrix
$\mathrm C^*$-subalgebra is the commutant of a unitary representation
of a finite group \cite[Theorem~4.1]{Schuch2010PEPS}. Their proof uses
\begin{align}
    \mathcal D\cong\bigoplus_i I_{d_i}\otimes M_{m_i}(\mathbb C),
    \label{eq:commutant_unital_blocks}
\end{align}
with no zero summand. Source traceability is
\path{Papers/1001.3807/paper_v3.tex}, lines 852--886. Earlier, however,
the source defines a $\mathrm C^*$-algebra simply as a linear matrix
space closed under multiplication and adjoint, at lines 561--563.
That wording alone does not require the ambient identity to belong to
$\mathcal D$.

The ambient identity belongs to every commutant. Consequently the
claim requires $I_D\in\mathcal D$. For example,
$\{\operatorname{diag}(z,0):z\in\mathbb C\}\subset M_2(\mathbb C)$
is closed under multiplication and adjoint, but cannot be a commutant
in $M_2(\mathbb C)$. This explains why the absence of a zero summand
in (\ref{eq:commutant_unital_blocks}) matters.

\section{The formal statement and the supported interpretation}
The library theorem is stated for a unital complex star subalgebra,
so membership of the ambient identity is part of the mathematical
object. It constructs the finite group and its unitary representation
from the subalgebra, and assumes no block decomposition or irreducible
representations as additional data. Its blueprint entry explicitly
states the unital qualification and is the unital version of
Theorem~4.1, rather than the unrestricted reading of the earlier
wording. This is a local correction to the source statement.

For a subalgebra whose identity is a proper projection $p$, the
appropriate ambient algebra is $pM_D(\mathbb C)p$, acting on
$p\mathbb C^D$. There the subalgebra is unital, and the corrected
commutant theorem applies. In the matrix-product-state application,
passing to this support removes unused bond directions. No claim
about the commutant in the original larger matrix algebra follows.

\bibliographystyle{plain}
\bibliography{references}
\end{document}
